\documentclass[12pt,a4paper]{report}
\usepackage[utf8]{inputenc}
\usepackage{amsmath}
\usepackage{amsfonts}
\usepackage{amssymb}
\usepackage{amsthm}
\usepackage{hyperref}

\usepackage{multicol}
\usepackage{fancyhdr}
\usepackage[inline]{enumitem}
\usepackage{tikz}
\usepackage{tikz-cd}
\usetikzlibrary{calc}
\usetikzlibrary{shapes.geometric}
\usetikzlibrary{positioning}
\usepackage[margin=0.5in]{geometry}
\usepackage{xcolor}

\hypersetup{
    colorlinks=true,
    linkcolor=blue,
    filecolor=magenta,      
    urlcolor=cyan,
    pdftitle={Tensors},
    pdfpagemode=FullScreen,
    }

%\urlstyle{same}

\newcommand{\CLASSNAME}{Math 5050 -- Special Topics: Tensor Analysis}
\newcommand{\STUDENTNAME}{Paul Carmody}
\newcommand{\ASSIGNMENT}{Chapter 10: Tensors on Surfaces }
\newcommand{\DUEDATE}{March 2026}
\newcommand{\PROFESSOR}{Professor Berchenko-Kogan}
\newcommand{\SEMESTER}{Spring 2026}
\newcommand{\SCHEDULE}{TBD}
\newcommand{\ROOM}{Remote}

\newcommand{\MMN}{M_{m\times n}}
\newcommand{\FF}{\mathcal{F}}

\pagestyle{fancy}
\fancyhf{}
\chead{ \fancyplain{}{\CLASSNAME} }
%\chead{ \fancyplain{}{\STUDENTNAME} }
\rhead{\thepage}

\fancyfoot{} % clear all footer fields
\fancyfoot[LE,RO]{\thepage}
\fancyfoot[LO,CE]{-------\\\ASSIGNMENT}
%\fancyfoot[CO,RE]{To: Dean A. Smith}
\input{../MyMacros}

\newcommand{\RED}[1]{\textcolor{red}{#1}}
\newcommand{\BLUE}[1]{\textcolor{blue}{#1}}

\newcommand{\SUPP}{\operatorname{supp}}
\newcommand{\CLOSURE}{\operatorname{closure of}}
\newcommand{\BD}{\operatorname{bd}}
\newcommand{\HHN}{\mathcal{H}^n}
\newcommand{\COFF}[3]{\Gamma^{#1}_{{#2}{#3}}}
\newcommand{\VK}[1]{\textbf{#1}}
\newcommand{\VKR}{\VK{R}}
\newcommand{\VKZ}{\VK{Z}}

\begin{document}

\begin{center}
	\Large{\CLASSNAME -- \SEMESTER} \\
	\large{ w/\PROFESSOR}
\end{center}
\begin{center}
	\STUDENTNAME \\
	\ASSIGNMENT -- \DUEDATE\\
\end{center} 

\noindent\textbf{Exercises: 10.1}
\begin{description}

\item \textbf{Exercise 213} Explain why it is not possible that $Z_\alpha^i Z_j^\alpha = \delta_j^i$.

\BLUE{First of all $Z_\alpha^i$ is a $3\times 2$ matrix and $Z_j^\alpha$ is a $2\times 3$ matrix, thus $Z_\alpha^i Z_j^\alpha$ is a $2 \times 2$ matrix when $\delta_j^i$ must be a $3 \times 3$ matrix.
}

\item \textbf{Exercise 214} Derive equation (10.15) from equation (10.14).

\section*{Section 10.6}
\item \textbf{Exercise 215} Show that $(\textbf{V}-\textbf{P})\cdot\textbf{N}$, where $\textbf{P}$ is defined by equation (10.42), is zero.
\begin{align*}
	\textbf{P} &= \PAREN{\textbf{V}\cdot\textbf{N}}\textbf{N} & (10.42)
\end{align*}

\BLUE{\begin{align*}
	(\textbf{V}-\textbf{P})\cdot\textbf{N} &= \PAREN{\textbf{V}-\PAREN{\textbf{V}\cdot\textbf{N}}\textbf{N}}\cdot\textbf{N} \\
	&= \textbf{V}\cdot\textbf{N} - (\textbf{V}\cdot\textbf{N})\textbf{N}\cdot\textbf{N} \\
	&= \textbf{V}\cdot\textbf{N}(1-\|N\|) \\
	&= 0
\end{align*}
}

\item \textbf{Exercise 216} Denote the tensor $N^iN_j$ by $P_j^i$.  Show that 
\begin{align*}
	P_j^i P_k^j &= P_k^i.
\end{align*}
\item \textbf{Exercise 217} To show $\textbf{V}-\textbf{T}$ that is orthogonal to the tangent plane, show that $(\textbf{V}-\textbf{T})\cdot \textbf{S}^\beta = 0$.
\item \textbf{Exercise 218} Denote the tensor $N^iN_j$ by $T_j^i$. Show that 
\begin{align*}
	T_j^iT_k^j=T_k^i,
\end{align*}

\BLUE{
\begin{align*}
		T_j^iT_k^j &= N^i(N_jN^j)N_k \\
		&= N^i\|N\|N_k \\
		&= N^iN_k
\end{align*}
}

\item\textbf{Exercise 219} Derive equation (10.41) from equation (10.55)
\item \textbf{Exercise 220} Show that the two terms in equation (10.74) are equal, thus
\begin{align*}
	N^iN_r &= \frac{1}{2}\delta_{rst}^{ijk}T_j^sT_k^i
\end{align*}According to equation (9.20), we find 
\begin{align*}
N^iN_r = \frac{1}{2}\PAREN{\delta_r^iT_j^jT_k^k-\delta_r^jT_j^iT_k^k+\delta_r^jT_j^kT_k^i-\delta_r^kT_j^jT_k^i+\delta_r^kT_j^iT_k^j-\delta_r^iT_j^kT_k^j}
\end{align*}which yields the following $\dots$.
\item \textbf{Exercise 221} Derive Equation (10.81) from the definition (10.80).
\item \textbf{Exercise 222} Derive Equaion (10.82).
\item \textbf{Exercise 223} Show that when the ambient space is referred to affine coordinates, the relationship between the surface and the ambient Christoffel symbols simplies to 
\begin{align*}
	\Gamma_{\beta\gamma}^\alpha = Z_i^\alpha\frac{\partial Z_\beta^i}{\partial S^\gamma}.
\end{align*}
\section*{Section 10.11}
\item \textbf{Exercise 224} Derive the objects given in this section.
\item \textbf{Exercise 225} Rederive the objects on the sphere by referring the ambient space to sperical coordinates $r_1, \theta_1, \phi_1$ in which the sphere is give by 
\begin{align*}
	r_1(\theta,\phi) &= R\\
	\theta_1(\theta, \phi) &= \theta \\
	\phi_1(\theta,\phi_) &= \phi
\end{align*}
\item \textbf{Exercise 226} Show that when the curve is referred to the arc length $s$, the above expression simplifies to the following
\begin{align*}
	N^i_, = \COLVECTOR{y'(x)\\x'(s)}; S_{\alpha\beta} =1 ; S^{\alpha\beta}=1; \sqrt{S}=1; \tilde{\Gamma} = 0
\end{align*}
\section*{Section 10.12}
\item \textbf{Exercise 228} Show that the when the curve is referred to the arc length $s$, the above expression simplify to the following
\begin{align*}
	Z_\alpha^i &= \COLVECTOR{r'(s)\\\theta'(s)}; N^i = \COLVECTOR{r(s)\theta'(s)\\-r(s)^{-1}r'(s)};\\
	S_{\alpha\beta}&=1; S^{\alpha\beta}=1\\
	\sqrt{S} &= 1; \tilde{\Gamma} = 0
\end{align*}
\end{description}
\end{document}
